\documentclass[../../../include/open-logic-section]{subfiles}

\begin{document}

\olfileid{sth}{spine}{rank}
\olsection{رتبه}

اوس چې پړاوونه مو د $V_\alpha$ په توګه تعريف کړي او پوهېږو چې هر سټ
د کوم پړاو فرعي سټ دے، د سټ \emph{رتبه} تعريفولے شو۔ په بداهتي توګه،
د $A$ رتبه هغه لومړۍ شېبه ده چې $A$ پکښې جوړېږي۔ په کره توګه:

\begin{defn}\ollabel{defnsetrank}
د هر سټ $A$ دپاره $\setrank{A}$ تر ټولو وړوکے هغه ترتيبي عدد $\alpha$
دے چې $A \subseteq V_\alpha$ دے۔
\end{defn}
\begin{prop}\ollabel{ranksexist}
	$\setrank{A}$ د هر $A$ دپاره شته۔
\end{prop}
\begin{proof}
	د تمرين په توګه پرېښودل شوے دے۔
\end{proof}
\begin{prob}
	د \olref[sth][spine][rank]{ranksexist} ثبوت ورکړئ۔
\end{prob}
\noindent
د رتبو ښه ترتيب موږ ته د ځينو مهمو پايلو د ثابتولو امکان راکوي:
\begin{prop}\ollabel{valphalowerrank}
د هر ترتيبي عدد $\alpha$ دپاره $V_\alpha = \Setabs{x}{\setrank{x} \in \alpha}$ دے۔
\end{prop}

\begin{proof}
که $\setrank{x} \in \alpha$ وي، نو $x \subseteq V_{\setrank{x}} \in
V_\alpha$ دے؛ نو $x \in V_\alpha$ دے، ځکه چې $V_\alpha$ توانمن دے
(دلته \olref[Valphabasic]{Valphabasicprops} څو ځلې کارول کېږي)۔ بالعکس،
که $x \in V_\alpha$ وي، نو $x \subseteq V_\alpha$ دے؛ له دې امله
$\setrank{x} \leq \alpha$ دے۔ اوس يو ساده نامتناهيت هاخوا استقرا ښيي چې
$x \notin V_{\setrank{x}}$ دے؛ نو د سټ رتبه له ورکړل شوي ترتيبي عدد
څخه کمه ده۔ (د سرچينې سمون، OLSTH-007: سرچينه په دې وروستۍ وينا کښې
د هماغه ورکړل شوي پړاو غړيتوب نفي کوي چې د ثبوت فرض ئې مني۔ د استقرا
د سمې پايلې پړاو د سټ خپله رتبه ده؛ دلته يوازې دغه شاخص سم شوے دے۔)
\end{proof}
\begin{prob}
	په \olref[sth][spine][rank]{valphalowerrank} کښې ياد شوے ساده
	نامتناهيت هاخوا استقرا بشپړ کړئ۔
\end{prob}

\begin{prop}\ollabel{rankmemberslower}
که $B \in A$ وي، نو $\setrank{B} \in \setrank{A}$ دے۔
\end{prop}

\begin{proof}
د \olref{valphalowerrank} له مخې
$A \subseteq V_{\setrank{A}} = \Setabs{x}{\setrank{x} \in
\setrank{A}}$ دے۔
\end{proof}
\noindent
په دې حقيقت داسې پايله ثابتولے شو چې د استقرا د يوې بڼې له لارې د
\emph{ټولو سټونو} په اړه د ويناوو د ثابتولو امکان راکوي:

\begin{thm}[د $\in$-استقرا شېما]
د هر فارمول $\phi$ دپاره:
\[
	\forall A((\forall x \in A)\phi(x) \lif \phi(A)) \lif \forall A \phi(A).
\]
\end{thm}

\begin{proof}
نقيض عکس به ثابتوو۔ نو فرض کړئ $\lnot \forall A
\phi(A)$ دے۔ د نامتناهيت هاخوا استقرا
(\olref[ordinals][basic]{ordinductionschema}) له مخې يو داسې سټ شته
چې $\phi$ نۀ دے او تر ټولو وړوکې ممکنه رتبه لري؛ يعنې داسې $A$ شته
چې $\lnot \phi(A)$ دے او $\forall x(\setrank{x} \in \setrank{A}
\lif \phi(x))$ دے۔ اوس که $x \in A$ وي، نو د
\olref{rankmemberslower} له مخې $\setrank{x} \in
\setrank{A}$ دے؛ له دې امله $\phi(x)$ دے؛ يعنې $(\forall x \in
A)\phi(x) \land \lnot \phi(A)$ دے۔
\end{proof}\noindent دا ځواکمنه پايله په غير رسمي توګه داسې بيانولے
شو: $\phi$ هله او يوازې هله \emph{وراثتي} دے چې، هر کله د يوه سټ
هر !!{element} د $\phi$ خاصيت ولري، هماغه سټ پخپله هم $\phi$ وي۔ نو
$\in$-استقرا دا درته وايي: که $\phi$ وراثتي وي، هر سټ $\phi$ دے۔

د رتبو بحث د اوس دپاره په دې توګه پاې ته رسوو چې هغه څو ادعاوې
ثابتې کړو چې څو ځلې مو ئې مخکښې اشاره کړې وه۔

\begin{prop}\ollabel{ranksupstrict}
$\setrank{A} = \supstrict_{x \in A}\setrank{x}$.
\end{prop}

\begin{proof}
$\alpha = \supstrict_{x \in A}\setrank{x}$ واخلئ۔ د
\olref{rankmemberslower} له مخې $\alpha \leq \setrank{A}$ دے۔ خو که
$x \in A$ وي، نو $\setrank{x} \in \alpha$ دے؛ نو د
\olref{valphalowerrank} له مخې $x \in V_\alpha$ دے، او له دې امله
$A \subseteq V_\alpha$ دے؛ يعنې $\setrank{A} \leq \alpha$ دے۔ نو
$\setrank{A} = \alpha$ دے۔
\end{proof}

\begin{cor}\ollabel{ordsetrankalpha}
د هر ترتيبي عدد $\alpha$ دپاره $\setrank{\alpha} = \alpha$ دے۔
\end{cor}

\begin{proof}
د نامتناهيت هاخوا استقرا دپاره فرض کړئ چې $\setrank{\beta} = \beta$
د هر $\beta \in \alpha$ دپاره دے۔ اوس د \olref{ranksupstrict} له مخې
$\setrank{\alpha} = \supstrict_{\beta \in
\alpha}\setrank{\beta} = \supstrict_{\beta \in \alpha}\beta = \alpha$
دے۔
%	First note that $\setrank{\alpha} \neq \beta$ for any $\beta \in
%	\alpha$. For otherwise we would have $\beta = \setrank{\beta} \in
%	\setrank{\alpha} = \beta$ by \olref{rankmemberslower}, i.e.,
%	$\beta \in \beta$, a contradiction. 
%
%	Now note that $\alpha \subseteq V_\alpha$. For $\alpha =
%	\Setabs{\beta}{\beta \in \alpha}$ by
%	\olref[sfr][ordinals][basic]{ordissetofsmallerord}. And if
%	$\beta \in \alpha$ then $\beta \subseteq V_{\setrank{\beta}} =
%	V_\beta \in V_\alpha$, hence $\beta \in V_\alpha$, by
%	\olref[sfr][spine][Valphabasic]{Valphabasicprops} twice. 
\end{proof}

په پاې کښې، د هغې پايلې يو لنډ ثبوت وړاندې کوو چې د
\olref[foundation]{sec} په پاې کښې مو ئې ژمنه کړې وه: $\ZFminus$ د
$\emph{منظموالی} \Rightarrow \emph{د بنسټ اصل}$ شرطي وينا ثابتوي۔
(ياد ولرئ چې په دې ثبوت کښې د ``رتبې'' مفهوم او
\olref{rankmemberslower} کارولے شو؛ ځکه لکه څنګه چې د دې برخې په
پيل کښې مو ياد کړل، دا د $\ZFminus + \text{منظموالی}$ له مخې وړاندې
کېدلے شي۔)

\begin{prop}[په $\ZFminus + \text{منظموالی}$ کښې کار کوو]\ollabel{zfminusregularityfoundation} د بنسټ اصل صدق کوي۔
\end{prop}

\begin{proof}
$A \neq \emptyset$ وټاکئ او داسې $B \in A$ واخلئ چې تر ټولو وړوکې
ممکنه رتبه لري۔ که $c \in B$ وي، نو د \olref{rankmemberslower} له مخې
$\setrank{c} \in \setrank{B}$ دے؛ نو $c \notin A$ دے، د $B$ د
انتخاب له امله۔
\end{proof}

\end{document}
